% GENERATED FILE. Edit canonical lessons, then run npm run generate:books.
% canonical-source-hash: fnv1a64:8efebe56ed4cbc2d
% canonical-lessons: ZH-C225-nongcun, ZH-C225-nongmin, ZH-C225-zhongdi, ZH-C225-zhai, ZH-C225-luying

\chapter{In the Country}
\label{ch:zh-in-the-country}
\hlchaptermodality{225}

\begin{chapteropening}
\textbf{By the end of this chapter:} \emph{I can talk about the countryside, farmers, working the land, picking fruit and camping.}

Complete the last lesson of chapter 225: I can talk about the countryside, farmers, working the land, picking fruit and camping.
\end{chapteropening}

\section[nóngcūn]{nóngcūn — the countryside; a farming village}

\label{lesson:ZH-C225-nongcun}



\begin{quote}
Before the new one: say the Mandarin for a butterfly, then the Mandarin for a bee.
\end{quote}

\subsection*{You'll want to know: nóngcūn}
\emph{nóngcūn} — \textquotedblleft{}the countryside; a farming village\textquotedblright{}.

\begin{quote}
\emph{wǒ lǎojiā zài nóngcūn} — My home town is out in the countryside.
\end{quote}

Its characters come later; for now it is a word for the ear.

\begin{cousinweb}
\emph{nóngcūn} is built from \emph{nóng} 'farming' and \emph{cūn} 'village', both inherited from Middle Chinese, \emph{nowng} and \emph{tshwon} (Guangyun, in Baxter's transcription), which glosses the first 'farming the fields' and the second 'a farmstead'.
\end{cousinweb}

\begin{grammarlens}[title={What you've built}]
One more word for saying what I do.
\end{grammarlens}

\subsection*{Your turn}
\begin{itemize}
\raggedright
  \item \emph{Say it:} \emph{nóngcūn}
  \item \emph{Say it:} \emph{nóngcūn}, once more
  \item \emph{Say it:} \emph{wǒ lǎojiā zài nóngcūn} — My home town is out in the countryside.
  \item \emph{Recall:} say the Mandarin for a butterfly, then the Mandarin for a bee, then say \emph{nóngcūn} again
\end{itemize}

\begin{tcolorbox}[breakable,colback=teal!4,colframe=teal!35!black,title={Before you move on}]
What does nóngcūn mean? (\textquotedblleft{}The countryside; a farming village\textquotedblright{}.) Say it once more.
\end{tcolorbox}
\section[nóngmín]{nóngmín — farmer}

\label{lesson:ZH-C225-nongmin}



\begin{quote}
Before the new one: say the Mandarin for a bee, then the Mandarin for the countryside.
\end{quote}

\subsection*{You'll want to know: nóngmín}
\emph{nóngmín} — \textquotedblleft{}farmer\textquotedblright{}.

\begin{quote}
\emph{tā bàba shì nóngmín} — His dad is a farmer.
\end{quote}

Its characters come later; for now it is a word for the ear.

\begin{cousinweb}
\emph{nóngmín} is built from \emph{nóng} 'farming' and \emph{mín} 'the people', both inherited from Middle Chinese, \emph{nowng} and \emph{mjin} (Guangyun, in Baxter's transcription), which explains the second, after an older book, as the multitude.
\end{cousinweb}

\begin{grammarlens}[title={What you've built}]
One more word for saying what I do.
\end{grammarlens}

\subsection*{Your turn}
\begin{itemize}
\raggedright
  \item \emph{Say it:} \emph{nóngmín}
  \item \emph{Say it:} \emph{nóngmín}, once more
  \item \emph{Say it:} \emph{tā bàba shì nóngmín} — His dad is a farmer.
  \item \emph{Recall:} say the Mandarin for a bee, then the Mandarin for the countryside, then say \emph{nóngmín} again
\end{itemize}

\begin{tcolorbox}[breakable,colback=teal!4,colframe=teal!35!black,title={Before you move on}]
What does nóngmín mean? (\textquotedblleft{}Farmer\textquotedblright{}.) Say it once more.
\end{tcolorbox}
\section[zhòngdì]{zhòngdì — to farm; to work the land}

\label{lesson:ZH-C225-zhongdi}



\begin{quote}
Before the new one: say the Mandarin for the countryside, then the Mandarin for a farmer.
\end{quote}

\subsection*{You'll want to know: zhòngdì}
\emph{zhòngdì} — \textquotedblleft{}to farm; to work the land\textquotedblright{}.

\begin{quote}
\emph{nóngmín měitiān zhòngdì} — Farmers work the land every day.
\end{quote}

Its characters come later; for now it is a word for the ear.

\begin{cousinweb}
\emph{zhòngdì} is built from \emph{zhòng} 'to plant' and \zh{地} 'land', both inherited from Middle Chinese, \emph{tsyowngH} and \emph{dijH} (Guangyun, in Baxter's transcription), which glosses the first 'to plant in the soil'.
\end{cousinweb}

\begin{grammarlens}[title={What you've built}]
One more word for saying what I do.
\end{grammarlens}

\subsection*{Your turn}
\begin{itemize}
\raggedright
  \item \emph{Say it:} \emph{zhòngdì}
  \item \emph{Say it:} \emph{zhòngdì}, once more
  \item \emph{Say it:} \emph{nóngmín měitiān zhòngdì} — Farmers work the land every day.
  \item \emph{Recall:} say the Mandarin for the countryside, then the Mandarin for a farmer, then say \emph{zhòngdì} again
\end{itemize}

\begin{tcolorbox}[breakable,colback=teal!4,colframe=teal!35!black,title={Before you move on}]
What does zhòngdì mean? (\textquotedblleft{}To farm; to work the land\textquotedblright{}.) Say it once more.
\end{tcolorbox}
\section[zhāi]{zhāi — to pick fruit; to pick flowers}

\label{lesson:ZH-C225-zhai}



\begin{quote}
Before the new one: say the Mandarin for a farmer, then the Mandarin for to work the land.
\end{quote}

\subsection*{You'll want to know: zhāi}
\emph{zhāi} — \textquotedblleft{}to pick fruit; to pick flowers\textquotedblright{}.

\begin{quote}
\emph{wǒmen qù nóngcūn zhāi píngguǒ} — We're going to the countryside to pick apples.
\end{quote}

Its characters come later; for now it is a word for the ear.

\begin{cousinweb}
\emph{zhāi} is inherited from Middle Chinese \emph{treak} (Guangyun, in Baxter's transcription), glossed 'to take with the hand'.
\end{cousinweb}

\begin{grammarlens}[title={What you've built}]
One more word for saying what I do.
\end{grammarlens}

\subsection*{Your turn}
\begin{itemize}
\raggedright
  \item \emph{Say it:} \emph{zhāi}
  \item \emph{Say it:} \emph{zhāi}, once more
  \item \emph{Say it:} \emph{wǒmen qù nóngcūn zhāi píngguǒ} — We're going to the countryside to pick apples.
  \item \emph{Recall:} say the Mandarin for a farmer, then the Mandarin for to work the land, then say \emph{zhāi} again
\end{itemize}

\begin{tcolorbox}[breakable,colback=teal!4,colframe=teal!35!black,title={Before you move on}]
What does zhāi mean? (\textquotedblleft{}To pick fruit; to pick flowers\textquotedblright{}.) Say it once more.
\end{tcolorbox}
\section[lùyíng]{lùyíng — to camp out}

\label{lesson:ZH-C225-luying}



\begin{quote}
Before the new one: say the Mandarin for to work the land, then the Mandarin for to pick fruit.
\end{quote}

\subsection*{You'll want to know: lùyíng}
\emph{lùyíng} — \textquotedblleft{}to camp out\textquotedblright{}.

\begin{quote}
\emph{zhōumò wǒmen qù lùyíng} — We're going camping at the weekend.
\end{quote}

Its characters come later; for now it is a word for the ear.

\begin{cousinweb}
\emph{lùyíng} is built from \emph{lù} 'dew, the open air' and \emph{yíng} 'a camp'. The second is inherited from Middle Chinese \emph{yweng} (Guangyun, in Baxter's transcription), glossed 'to build'; the dictionary explains the first, \emph{luH} there, as dew. How the pair came to mean camping is not settled.
\end{cousinweb}

\begin{grammarlens}[title={What you've built}]
One more word for saying what I do.
\end{grammarlens}

\subsection*{Your turn}
\begin{itemize}
\raggedright
  \item \emph{Say it:} \emph{lùyíng}
  \item \emph{Say it:} \emph{lùyíng}, once more
  \item \emph{Say it:} \emph{zhōumò wǒmen qù lùyíng} — We're going camping at the weekend.
  \item \emph{Recall:} say the Mandarin for to work the land, then the Mandarin for to pick fruit, then say \emph{lùyíng} again
\end{itemize}

\begin{tcolorbox}[breakable,colback=teal!4,colframe=teal!35!black,title={Before you move on}]
What does lùyíng mean? (\textquotedblleft{}To camp out\textquotedblright{}.) Say it once more.
\end{tcolorbox}
